\section{Benchmark and evaluation}
\label{sec:methods}

\paragraph{Datasets and chemical partitions.}
We use all 22 TDC ADMET datasets as separate tasks: 13 binary and nine regression. Molecules are
standardized to parent structures with the ChEMBL pipeline \citep{bento2020chembl} and keyed by
parent Standard InChIKey. We keep repeated measurements but assign all records for one compound to
the same fold. Inputs concatenate 2048-dimensional radius-2 Morgan count fingerprints
\citep{rogers2010morgan} and 200 CDF-normalized RDKit2D descriptors computed with DescriptaStorus
2.8.0 \citep{kelley2024descriptastorus}. Every pipeline is fitted on the log target for four endpoints:
\dataset{vdss\_lombardo}, \dataset{half\_life\_obach},
\dataset{clearance\_microsome\_az}, and \dataset{clearance\_hepatocyte\_az}. Predictions are
transformed back before scoring.

Five outer folds are formed from Taylor--Butina clusters at Tanimoto distance cutoff 0.5
\citep{butina1999clustering}. For each outer fold, we recompute the clusters within the outer-training
compounds and divide them into a dedicated calibration set and a model-development set. Complete
clusters are added to the calibration set until it contains approximately 20\% of the observations.
Only the model-development set enters five-fold inner cross-validation. We average each
configuration's five TDC-score trajectories step by step, choose its best step on that mean curve,
and select the configuration with the best resulting score. We refit it on the full model-development
set and fit the post-hoc calibrator on the untouched calibration set.

\paragraph{Prediction pipelines.}
We compare LightGBM \citep{ke2017lightgbm} with five neural pipelines: a deterministic multilayer
perceptron (MLP), MC dropout \citep{gal2016dropout}, checkpoint-resumed stochastic
weight averaging Gaussian (SWAG) \citep{maddox2019swag}, last-layer Laplace
\citep{daxberger2021laplace}, and a post-hoc variational Bayesian last layer (VBLL)
\citep{harrison2024vbll}. Their ten-member
counterparts are deep ensemble \citep{lakshminarayanan2017deep}, ensemble MC dropout
\citep{filos2019systematic}, MultiSWAG \citep{wilson2020bayesian}, mixture of Laplace approximations
(MoLA) \citep{eschenhagen2021mixtures}, and ensemble post-hoc VBLL, giving eleven pipelines. The ten
neural pipelines use the same MLP architecture, while LightGBM provides the alternative model family;
all eleven use the same fingerprint-and-descriptor input.

LightGBM and the MLP backbones use binary cross-entropy for the 13 classification tasks, L1 loss for
the five MAE-scored regression tasks, and L2 loss for the four Spearman-scored regression tasks.
The post-hoc VBLL layer is fitted by minimizing the negative ELBO.

LightGBM and the deterministic MLP are selected from 32 presampled Sobol configurations using the
mean inner-fold TDC metric. All neural UQ pipelines inherit the MLP architecture, hyperparameters,
and training epoch selected for deterministic point prediction. MC dropout, SWAG, and Laplace use
the selected checkpoint; post-hoc VBLL fits a variational last layer over its frozen features. For
each feature map, we select among 81 last layers spanning nine prior and nine Wishart scales by the
final ELBO on the model-development rows; ensemble post-hoc VBLL repeats this for every member.
Ensembles combine ten independently initialized refits of the same selection. This holds backbone
selection fixed across neural UQ constructions. Fits run 2,000 LightGBM rounds or 300 MLP epochs,
with the deployed step selected from the inner validation trajectory.

Deterministic MLP and MC dropout use the same selected network and training-time dropout. At
inference, deterministic MLP disables dropout and uses one forward pass, whereas MC dropout keeps
dropout active and averages 30 stochastic passes. Starting from the selected checkpoint, SWAG runs
ten epochs of stochastic gradient descent with momentum and collects up to four parameter snapshots
per epoch from the first epoch. The resulting Gaussian approximation has diagonal-plus-rank-20 covariance
and is evaluated with 30 posterior weight draws. MultiSWAG
uses ten independently refitted SWAG members with equal weight, each evaluated with 30 posterior
weight draws. For regression, its combined variance includes both the average uncertainty within
members and the variance between member means. For binary prediction, we average the ten member
probabilities after averaging each member's 30 draws.

\paragraph{Calibration and evaluation measures.}
We estimate post-hoc calibration parameters from the dedicated calibration set. For binary tasks,
we apply temperature scaling to the logit of each pipeline's final marginal probability after
averaging ensemble members \citep{guo2017calibration}. We evaluate the resulting probabilities with
Brier score, a proper score that reflects both calibration and refinement.

Regression requires different calibration procedures because the pipelines supply different forms
of uncertainty. For LightGBM and the deterministic MLP, we estimate one constant residual standard
deviation. For MC dropout, SWAG, deep ensemble, ensemble MC dropout, and MultiSWAG, we estimate an
additional observation-noise variance and add it to the supplied predictive variance. For Laplace,
post-hoc VBLL, MoLA, and ensemble post-hoc VBLL, we multiply the predictive standard deviation by a
fitted non-negative factor. These adjusted distributions define central 90\% intervals, evaluated
by interval score. For the four log-trained endpoints, we estimate calibration parameters and form
intervals on the log target, then exponentiate the interval limits and score them in the endpoint's
original units.

Our main measures are the endpoint-specific TDC metric, Brier score for binary tasks, and interval
score for regression tasks. Repeated labels and predictions are averaged within compound for the
TDC metrics; Brier and interval scores use individual observations. The appendix also reports
AUROC, average precision, NLL, and ECE for binary tasks, and MAE, RMSE, Spearman correlation,
Gaussian NLL, and coverage error for regression.

\paragraph{Full grid and comparison across datasets.}
For each dataset and method under each evaluation measure, we average exactly
$10\times5\times5=250$ outer-fold scores: ten separately generated chemical partitions, five model seeds,
and five outer folds. These 250 values estimate one score. We rank the eleven methods within each
dataset, with rank 1 best and average ranks for exact score ties. Each dataset contributes one score
and one rank to the cross-dataset comparison. The 1,826,000 from-scratch model fits are itemized in
Appendix~Table~\ref{tab:compute-counts}.

For each evaluation measure, the tie-corrected Friedman statistic with the Iman--Davenport $F$
approximation tests whether method ranks differ overall across datasets \citep{demsar2006statistical}.
We then compare all 55 method pairs using two-sided normal tests of their mean-rank differences and
apply Holm's procedure to this fixed all-pairs family, controlling the family-wise error rate at 5\%
\citep{garcia2008extension}. We report the comparisons involving the method with the best average
rank. A non-significant comparison means that the data do not rule out the alternative as best; it
does not establish equal performance.

\paragraph{Per-dataset comparisons.}
Average ranks do not show what can be resolved on one endpoint, so we also compare methods on each
dataset's original score scale. We average the five model seeds and five outer folds within each
chemical partition, giving 10 paired complete-partition scores per method. A repeated-measures
analysis of variance first tests whether any method differs; if so, Hsu's MCB procedure at the 5\%
level assesses which methods can be ruled out as best \citep{hsu1984constrained}. We call the
observed mean in the favorable direction the empirical best. Appendix~\ref{app:dataset-scores}
gives the decision threshold and its coverage interpretation.

The ten partitions are generated separately but reuse the same finite set of compounds; their
variation therefore measures sensitivity to the chemical partition rather than population-sampling
uncertainty for new compounds. The analysis also assumes normally distributed errors and one residual
variance across methods. Its $p$-values should therefore be read as conditional on these assumptions,
the observed datasets, and the completed evaluation grid.

\paragraph{Stability under smaller evaluation designs.}
We use finite-grid subsampling to ask how much of the completed evaluation is needed to recover its
conclusions. We compare one held-out fold with one or more complete five-fold CVs, using either one
model seed or the average of all five. Partition and fold levels are selected without replacement
and shared across methods and datasets. Every reported design enumerates all distinct paired
subsets. For each, we recalculate
the empirical winners, mean ranks, Holm decisions, and, when at least two blocks are available, the
per-dataset MCB conclusions. The resulting rates measure agreement with the completed grid.
